% =====================================================
% Rolle in Physik und Technik
% =====================================================
\subsection{Rolle in Physik und Technik}
\label{sec:rolle-physik-technik}

Das Elektron\index{Elektron} ist nicht nur ein Gegenstand der Grundlagenphysik\index{Grundlagenphysik}.
Es ist zugleich eines der wichtigsten Teilchen der modernen Technik. Kaum ein anderer
Baustein verbindet so viele Bereiche: Elektrotechnik\index{Elektrotechnik},
Atomphysik\index{Atomphysik}, Chemie\index{Chemie}, Festkörperphysik\index{Festkörperphysik},
Elektronik\index{Elektronik}, Messtechnik\index{Messtechnik}, Kommunikationstechnik\index{Kommunikationstechnik}
und moderne Informationsverarbeitung\index{Informationsverarbeitung}.

In der klassischen Elektrotechnik erscheint das Elektron zunächst indirekt. Man misst
Spannung\index{Spannung}, Stromstärke\index{Stromstärke}, Widerstand\index{Widerstand},
Leistung\index{Leistung} und Magnetfelder\index{Magnetfeld}, ohne jedes einzelne Elektron
zu verfolgen. Dennoch beruhen elektrische Ströme in metallischen Leitern auf der
geordneten Bewegung vieler Elektronen. Aus der mikroskopischen Bewegung einzelner
Ladungsträger\index{Ladungsträger} entsteht eine makroskopisch messbare Größe:
der elektrische Strom\index{elektrischer Strom}.

In der Atomphysik bestimmt das Elektron den Aufbau der Atomhülle\index{Atomhülle}.
Die Verteilung der Elektronen um den Atomkern\index{Atomkern} entscheidet darüber,
welche Energieniveaus\index{Energieniveau} ein Atom besitzt, welche Spektrallinien\index{Spektrallinie}
es aussendet oder absorbiert und wie es auf äußere Felder reagiert. Damit steht das
Elektron im Zentrum des Verständnisses von Atomen, Lichtemission\index{Lichtemission}
und Spektroskopie\index{Spektroskopie}.

In der Chemie ist das Elektron der eigentliche Träger der Bindung. Chemische
Bindungen\index{chemische Bindung} entstehen durch die Umordnung und gemeinsame Nutzung
von Elektronen. Ob ein Stoff reaktionsfreudig oder stabil ist, ob er leitfähig,
farbig, magnetisch oder isolierend wirkt, hängt wesentlich von der Struktur seiner
Elektronenhülle ab. Die Chemie ist daher in einem tiefen Sinn eine Physik der
Elektronen in Atomen und Molekülen.

In Festkörpern\index{Festkörper} zeigt sich die Bedeutung des Elektrons besonders
deutlich. Ob ein Material ein Leiter\index{Leiter}, Halbleiter\index{Halbleiter} oder
Isolator\index{Isolator} ist, hängt von den möglichen Bewegungszuständen seiner
Elektronen ab. Die Bandstruktur\index{Bandstruktur} eines Materials entscheidet darüber,
ob Elektronen frei beweglich sind oder an bestimmte Zustände gebunden bleiben. Diese
Einsicht bildet die Grundlage der modernen Halbleitertechnik\index{Halbleitertechnik}.

Damit führt das Elektron direkt in die moderne Elektronik. Dioden\index{Diode},
Transistoren\index{Transistor}, integrierte Schaltungen\index{integrierte Schaltung},
Sensoren\index{Sensor}, Speicherchips\index{Speicherchip}, Computer\index{Computer}
und Kommunikationssysteme beruhen darauf, dass Elektronen gezielt gesteuert werden
können. Die digitale Welt ist ohne Elektronen nicht denkbar: Jedes Schalten, Speichern,
Messen und Übertragen von Information nutzt letztlich elektrische Ladungen und ihre
kontrollierte Bewegung.

Auch in der modernen Physik bleibt das Elektron ein Schlüsselteilchen. Es ist ein
Prüfstein für die Quantenmechanik\index{Quantenmechanik}, für die Relativitätstheorie\index{Relativitätstheorie}
und für die Quantenelektrodynamik\index{Quantenelektrodynamik}. Präzisionsmessungen
am Elektron gehören zu den genauesten Tests physikalischer Theorien. Gerade weil das
Elektron so einfach erscheint, eignet es sich hervorragend, um die Grenzen unseres
Verständnisses auszuloten.

Die Rolle des Elektrons reicht also von der Steckdose bis zum Teilchenbeschleuniger,
vom chemischen Molekül bis zum Mikrochip, vom Magnetfeld bis zur Quantenfeldtheorie.
Es ist zugleich Arbeitsteilchen der Technik und Testobjekt der Grundlagenphysik.
Genau diese Doppelrolle macht es zu einem idealen Ausgangspunkt für dieses Buch.

\begin{PhysicsBox}[Vom Alltag zur Grundlagenphysik]
	\label{box:vomAlltag}
	Das Elektron verbindet Alltagsphänomene mit moderner Theorie. Es erklärt den
	elektrischen Strom in Leitern, die chemische Bindung in Molekülen, die Eigenschaften
	von Halbleitern und die Funktionsweise elektronischer Bauteile.
	
	Gleichzeitig ist es ein zentrales Objekt der Quantenmechanik, der relativistischen
	Physik und der Quantenelektrodynamik. Das Elektron steht damit an der Schnittstelle
	von Technik, Materie und moderner Grundlagenphysik.
\end{PhysicsBox}